% ================================================================
% Section: Model Compression Techniques for \glspl{slm}
% Merged from main (2).tex — Section IV + sections/background.tex
% ================================================================

\section{Model Compression Techniques for SLMs}
\label{sec:compression}

Model compression is central to the \gls{slm} ecosystem: it is the primary pathway for reducing large pre-trained models to sizes deployable on edge hardware, and it provides a complementary route to training-from-scratch. The principal techniques are knowledge distillation, quantization, and pruning; in practice, compound pipelines combining all three have become the state of the art.

\subsection{How Compression Results Are Reported Here}
\label{subsec:compression_reporting}

Comparisons of compression techniques are routinely presented as a single
``compression ratio'' set against a single ``accuracy retention'' percentage. An
earlier version of this survey did the same, and the presentation does not
survive scrutiny, so we set out here why and what replaces it.

\emph{There is no single compression ratio.} Quantization reduces bits per
weight, leaving the parameter count untouched; pruning reduces the parameter
count, and reduces memory only where the sparsity pattern is one the runtime can
actually exploit; distillation reduces the parameter count of the student but is
better understood as reducing the training compute needed to reach a given
quality, since a distilled 4B model and a from-scratch 4B model have the same
deployment cost. Collapsing these into one number equates quantities that trade
off against different budgets. Distillation in particular has no well-defined
universal ratio at all: the student's size is chosen by the practitioner, so
``$10$--$100\times$'' describes a range of design choices in the literature, not a
property of the method.

\emph{There is no baseline-free accuracy retention.} A statement that 4-bit
quantization retains ``about 96\%'' of accuracy is uninterpretable without the
baseline precision, the model, the model's size, the task, the metric, and the
rule by which per-task retentions were aggregated. The dependence on each of
these is large: retention under \gls{ptq} falls with model size
\citep{slmquant2025}, differs by task family, with mathematical and multi-step
reasoning degrading faster than knowledge recall, and differs by method at
matched bit-width \citep{lin2024awq}. A single band averages across all of that
variation and then presents the average with two significant figures.

We therefore replace the single table with two. Table~\ref{tab:compression_axes}
separates the reductions each technique delivers along five distinct axes, which
are not interchangeable and are not summed. Table~\ref{tab:compression_quality}
reports what the reviewed literature establishes about quality cost, stated as a
direction and a scope condition with the source attached, rather than as a
number without a baseline. Where the literature supports a numeric claim under
stated conditions, the number appears in the text with its conditions; where it
does not, no number is given, and Section~\ref{subsec:challenges_verification}
specifies the experiment that would supply one.

\begin{table}[t]
\caption{What each compression technique reduces, by axis. Entries are what the technique does by construction, not measurements. ``Parameters'' counts stored weights; ``memory'' is realized weight footprint on a runtime that can exploit the format; ``training compute'' is the cost of producing the compressed artifact relative to training an equivalent model from scratch; ``latency'' is decode-time speedup on memory-bandwidth-bound hardware. A dash means the technique does not act on that axis.}
\label{tab:compression_axes}
\centering
\footnotesize\adjustbox{max width=\linewidth}{%
\begin{tabular}{@{}lp{1.7cm}p{1.9cm}p{2.1cm}p{2.4cm}p{2.1cm}@{}}
\toprule
\textbf{Technique} & \textbf{Parameter reduction} & \textbf{Weight-memory reduction} & \textbf{Training-compute reduction} & \textbf{Decode latency effect} & \textbf{Retraining required} \\
\midrule
PTQ, 8-bit          & ---            & $2\times$ vs.\ FP16 & ---  & Proportional to bandwidth saved & No \\
PTQ, 4-bit          & ---            & $3.2$--$3.4\times$ realized (Table~\ref{tab:memory_footprint}) & --- & Proportional to bandwidth saved, batch~1 & No \\
PTQ, 2-bit          & ---            & ${\sim}6\times$ realized & --- & As above & No; \gls{qat} in practice \\
Unstructured pruning, 50\% & $2\times$ & Only on sparse-capable runtimes & --- & None on dense kernels & Calibration only \\
Structured pruning  & Set by target width/depth & Proportional to parameters removed & --- & Proportional to parameters removed & Yes, to recover quality \\
Knowledge distillation & Set by student size & Proportional to student size & Large; the student is trained rather than the teacher & Proportional to student size & Yes, student training \\
Prune + distill + quantize & Set by target & Product of the pruning and quantization factors & $40\times$ reported for Minitron \citep{minitron2024} & Product of the above & Yes \\
\bottomrule
\end{tabular}%
}
\end{table}

\begin{table}[t]
\caption{What the reviewed literature establishes about the quality cost of each technique in the sub-10B regime, with the scope condition under which it was established. Claims are graded by evidence class (Section~\ref{subsec:reported_grading}): \evB~independent evaluation, \evC~vendor- or author-reported, \evD~synthesis across sources by the present authors.}
\label{tab:compression_quality}
\centering
\footnotesize\adjustbox{max width=\linewidth}{%
\begin{tabular}{@{}lp{6.4cm}p{3.2cm}c@{}}
\toprule
\textbf{Technique} & \textbf{What the literature supports} & \textbf{Scope of the evidence} & \textbf{Grade} \\
\midrule
PTQ, 8-bit & Effectively lossless on knowledge and mathematical benchmarks; no consistent degradation reported & Widely reproduced across model families & \evB \\
PTQ, 4-bit & A small but real accuracy cost that grows as model size falls, and is larger on multi-step reasoning than on knowledge recall & Reported for the sub-4B band \citep{slmquant2025} & \evC \\
PTQ, 4-bit, method comparison & \gls{awq} retains more accuracy than GPTQ at matched bit-width by protecting activation-salient channels & Author-reported, with independent reuse \citep{lin2024awq} & \evC \\
PTQ, 2-bit & \gls{ptq} degrades unacceptably; \gls{qat} recovers a meaningful part of the loss, and the margin over \gls{ptq} widens as bit-width falls & Vendor pipelines with full training-stack control \citep{gunter2025afm} & \evC \\
Unstructured pruning, 50\% & Perplexity increase below 1\% at 175B scale in one pass \citep{frantar2023sparsegpt}; the tolerable sparsity falls sharply as the base model shrinks & Established at large scale; extrapolation to sub-4B is not supported & \evC \\
Quantization vs.\ pruning & Quantization preserves fidelity, multilingual perplexity, and reasoning accuracy better than pruning at matched compression; \gls{awq} 4-bit beats 50\% Wanda on 6 of 7 tasks & Controlled study: 6 models, 0.5--3.8B, 7 languages, 7 tasks \citep{kurz2025pruning_vs_quant} & \evB \\
Prune + distill + quantize & A 4B Minitron matches from-scratch 4B models on commonsense, code, and math at $40\times$ less training compute & Single pipeline, single teacher \citep{minitron2024} & \evC \\
Domain compound pipeline & A clinical model under 1B retains reported accuracy within 5\% of its 7B counterpart on ICD coding and discharge summarization & Single study, single domain \citep{hasan2024optimclm} & \evC \\
Fairness under compression & Degradation is uneven across languages and dialects, and pruning is worse than quantization on multilingual tasks & Inferred across two sources \citep{blodgett2020language,kurz2025pruning_vs_quant} & \evD \\
\bottomrule
\end{tabular}%
}
\end{table}

% ----------------------------------------------------------------
\subsection{Knowledge Distillation}
\label{subsec:compression_kd}

\gls{kd} \citep{hinton2015distilling} transfers the ``dark knowledge'' embedded in a teacher model's soft output distribution to a smaller student model. The student is trained to minimize a combination of the standard cross-entropy loss on labeled data and a Kullback--Leibler (KL) divergence loss with respect to the teacher's output logits at a temperature $T > 1$. This approach was pioneered for BERT compression (DistilBERT \citep{sanh2019distilbert}, TinyBERT \citep{jiao2020tinybert}) and has since been substantially extended for generative \glspl{slm}.

\subsubsection{White-Box vs.\ Black-Box KD}

The survey by \citet{yang2024kd_survey} classifies \gls{llm} distillation into two paradigms:

\emph{White-box \gls{kd}} has access to the teacher's intermediate representations. TinyBERT \citep{jiao2020tinybert} pioneered layer-by-layer feature alignment, matching student hidden states, attention matrices, and embedding layers to teacher counterparts via MSE loss. In the decoder-only \gls{llm} setting, this requires pairing student and teacher layers, which is non-trivial when sizes differ substantially, and is computationally expensive but yields superior retention of teacher capabilities.

\emph{Black-box \gls{kd}} uses only the teacher's generated outputs, either token probabilities or free-form text, as the supervision signal. Orca \citep{ho2022orca} exemplifies text-output black-box \gls{kd}: GPT-4 generates step-by-step reasoning traces, which a LLaMA-based student learns to replicate. DeepSeek-R1 \citep{deepseek2025r1} extends this to long chain-of-thought reasoning distillation, producing 7B and 14B students that dramatically outperform identically-sized from-scratch models on mathematical reasoning.

\subsubsection{Chain-of-Thought Distillation}

A particularly productive form of black-box \gls{kd} is \gls{cot} distillation \citep{wei2022cot}: the teacher model generates explicit reasoning steps, and the student learns to produce or internalize these steps. \citet{feng2024counterfactual} demonstrate counterfactual distillation, which augments \gls{cot} traces with contrastive incorrect reasoning, as a way to sharpen the student's reasoning boundaries. \citet{liao2024nesycd} propose Neural-Symbolic Collaborative Distillation, using a neuro-symbolic teacher to provide structured reasoning supervision that is more transferable to small students than purely neural \gls{cot}.

\subsubsection{Data Augmentation in KD}

The survey by \citet{xu2024kd_survey} highlights data augmentation (DA) as a central paradigm within \gls{kd} frameworks. When the student lacks sufficient coverage of the teacher's knowledge space, the teacher can generate synthetic training instances targeting specific skill gaps, the mechanism underlying the Phi family's success \citep{gunasekar2023phi1}. The efficient \gls{kd} approach of \citet{ballout2024efficient_kd} demonstrates that targeted data augmentation can reduce the student training compute by 60\% while preserving 95\% of the teacher's task-specific accuracy.

\subsubsection{Self-Distillation}

An emerging approach eliminates the external teacher entirely. \citet{fu2024dynamic_selfd} propose Dynamic Self-Distillation, in which the model uses its own predictions from previous mini-batches as soft targets during fine-tuning, providing a form of implicit regularization that reduces overfitting on small fine-tuning datasets, a critical challenge for \glspl{slm}.

% ----------------------------------------------------------------
\subsection{Quantization}
\label{subsec:compression_quantization}

Quantization reduces the numerical precision of model weights and/or activations from 32-bit floats (FP32) or 16-bit (FP16/BF16) to integers or low-bit floats (INT8, INT4, FP4, NF4, INT2). This directly reduces memory footprint and, on hardware with native integer arithmetic (\glspl{npu}, modern CPUs), reduces inference latency.

\subsubsection{Post-Training Quantization (PTQ)}
\glsunset{ptq}

PTQ is performed after training is complete, requiring only a small calibration dataset (typically 128--512 samples) and no gradient computation. It is the dominant quantization paradigm because it is cheap and applies to any model without access to training infrastructure.

GPTQ \citep{frantar2023gptq} applies second-order weight compression: it uses an approximation of the inverse Hessian of the loss with respect to each layer's weights to determine optimal rounding decisions, minimizing the layer-wise squared error introduced by quantization. For \glspl{slm}, GPTQ 4-bit quantization of a 7B model reduces \gls{vram} from ${\sim}14$~GB to ${\sim}4$~GB.

\gls{awq} \citep{lin2024awq} observes that not all weights are equally important for activation output: 1\% of weights (those corresponding to large-magnitude activation channels) account for disproportionate reconstruction error. \gls{awq} identifies these ``salient'' channels using a small calibration set and applies per-channel scaling to protect them before quantization, consistently outperforming GPTQ on complex reasoning tasks and winning the MLSys 2024 Best Paper award.

\gls{gguf} is a quantization file format popularized by \texttt{llama.cpp} \citep{gerganov2023llamacpp} that packages model weights, vocabulary, and architecture metadata in a single binary file. It supports a range of quantization levels (Q2\_K through Q8\_0), letting users trade off quality against memory footprint along a near-continuous spectrum, and it is the format in which most sub-4B models are actually distributed for CPU-only laptop and single-board deployment.

SqueezeLLM \citep{kim2024squeezellm} introduced sensitivity-based non-uniform quantization with dense-and-sparse decomposition, showing that sensitivity varies dramatically across weight dimensions and enabling higher accuracy at ultra-low bit-widths.

\subsubsection{Quantization-Aware Training (QAT)}
\glsunset{qat}

QAT integrates simulated quantization into the training loop, allowing the model to learn weight values that are robust to quantization error. This is particularly valuable for aggressive quantization (2-bit, 3-bit) where \gls{ptq} accuracy drops unacceptably. \gls{qat} recovers a meaningful part of the accuracy that \gls{ptq} loses at the same bit-width, and the margin widens as bit-width falls, but it requires access to training data and significant compute, which is why it appears mainly in vendor pipelines with control of the full training stack, such as Apple's 2-bit on-device quantization \citep{gunter2025afm}.

\subsubsection{Quantization versus Pruning in the Sub-10B Regime}

\citet{kurz2025pruning_vs_quant} report a systematic comparison of \gls{ptq} quantization (GPTQ, \gls{awq}) against post-training pruning (SparseGPT, Wanda) specifically in the \gls{slm} regime (0.5B--3.8B parameters, six models, seven languages, seven downstream tasks). Their primary conclusion: \emph{quantization consistently outperforms pruning in preserving model fidelity, multilingual perplexity, and reasoning accuracy}. \gls{awq} at 4-bit outperforms 50\% Wanda pruning on 6 of 7 tasks. Quantization's advantage is most pronounced on language modeling and multilingual tasks; the gap narrows on complex knowledge tasks, where both methods show similar degradation under high compression.

% ----------------------------------------------------------------
\subsection{Pruning}
\label{subsec:compression_pruning}

Pruning removes parameters from a pre-trained model, either individually (unstructured) or in structured groups (entire heads, layers, or channels). Structured pruning yields hardware-friendly sparse models compatible with standard matrix multiplication libraries; unstructured pruning achieves higher compression ratios but requires special sparse hardware or software to realize latency benefits.

\subsubsection{Unstructured Pruning}

SparseGPT \citep{frantar2023sparsegpt} extends the GPTQ framework to weight removal: it prunes weights layer-by-layer using inverse-Hessian information to minimize the reconstruction error introduced by zeroing each weight. It can prune the 175B GPT-3 to 50\% sparsity in a single forward pass without weight updates, with a perplexity increase below 1\%.

Pruning by weights and activations (Wanda) \citep{sun2024wanda} proposes a simpler criterion: remove the weights with the smallest product of weight magnitude and corresponding input-activation norm. This needs only a forward pass over the calibration data (no second-order information) and runs about $10\times$ faster than SparseGPT while matching its accuracy at 50\% sparsity.

\subsubsection{Structured Pruning}

Structured pruning removes entire attention heads, MLP channels, or transformer layers, yielding immediately deployable models without special sparsity libraries. Layer pruning, which removes the least important layers as measured by cosine similarity between adjacent layer outputs, can remove a modest fraction of a model's depth with limited benchmark degradation, though the tolerable fraction falls sharply as the base model gets smaller. The Minitron pipeline \citep{minitron2024} applies depth and width pruning simultaneously, guided by importance scores derived from activation magnitudes over a calibration set, and is the best-documented instance of the approach at \gls{slm} scale.

% ----------------------------------------------------------------
\subsection{Combined Compression Pipelines}
\label{subsec:compression_combined}

The state of the art in \gls{slm} compression combines pruning, distillation, and quantization in a multi-stage pipeline. NVIDIA Minitron \citep{minitron2024} achieves a $40\times$ reduction in training compute relative to training Llama 3.1 8B from scratch by: (1)~structured pruning to target size; (2)~knowledge distillation from the original 8B teacher to the pruned student; (3)~post-training quantization for deployment. The resulting 4B Minitron model matches the performance of 4B models trained from scratch on commonsense, code, and math benchmarks.

OptimCLM \citep{hasan2024optimclm} applies a similar pipeline to clinical language models, combining all three compression stages to produce models under 1B that achieve within 5\% of clinical \gls{nlp} accuracy of their 7B counterparts. These results suggest that the combination of pruning (for architectural reduction) + distillation (for capability recovery) + quantization (for memory reduction) represents the practical frontier of \gls{slm} compression.
